% =====================================================================
% chapters/05-discussion.tex
% PT: Capítulo 5 -- Discussão. Exemplo de: referências cruzadas entre
%     capítulos (várias de uma vez), variantes dos acrónimos (\acrshort,
%     \acrlong, \acrfull), lista itemize e nota de rodapé.
% EN: Chapter 5 -- Discussion. Shows: cross-references between chapters
%     (several at once), acronym variants (\acrshort, \acrlong, \acrfull),
%     an itemize list and a footnote.
% =====================================================================
\chapter{Discussion}
\label{chap:discussion}

% PT: Parágrafo de abertura do capítulo.
% EN: Chapter opening paragraph.
This chapter shows how to refer back to earlier material. In a real
dissertation, the discussion interprets the results, relates them to the
literature, answers the research questions and states the limitations of the
work.

% ---------------------------------------------------------------------
% PT: Exemplo em inglês
% EN: English example
% ---------------------------------------------------------------------
\section{English example}
\label{sec:discussion:en}

\subsection{Interpretation of the Results}
\label{sec:discussion:interpretation}

% PT: \cref aceita várias etiquetas separadas por vírgulas e agrupa-as:
%     \cref{fig:a,fig:b} -> "Figures 3.1 and 4.1".
% EN: \cref accepts several labels separated by commas and groups them:
%     \cref{fig:a,fig:b} -> "Figures 3.1 and 4.1".
The results in \cref{tab:results} follow the pipeline of \cref{fig:pipeline},
and \cref{fig:pipeline,fig:example-file} can be cited together. The loss in
\cref{eq:loss} was the only training objective, as described in
\cref{sec:method:model}. These references stay correct when chapters,
figures or tables are added or reordered, because the numbers are generated
by \LaTeX.

% PT: Variantes dos acrónimos.
% EN: Acronym variants.
The acronym commands can also be used directly: \acrshort{cnn} prints the
short form, \acrlong{cnn} prints the long form, and \acrfull{ml} always
prints both, regardless of earlier uses.

\subsection{Limitations}
\label{sec:discussion:limitations}

Every study has limitations, and stating them makes the conclusions more
credible. Common limitations include:

\begin{itemize}
  \item the size and representativeness of the data;
  \item the number of methods and settings compared;
  \item the hardware and time available for the experiments.\footnote{State
    the limitations that actually affect your conclusions, not a generic
    list.}
\end{itemize}

% ---------------------------------------------------------------------
% PT: Exemplo em português
% EN: Portuguese example
% ---------------------------------------------------------------------
\section{Exemplo em português}
\label{sec:discussion:pt}

\begin{otherlanguage}{portuguese}

Os resultados da Tabela~\ref{tab:results} devem ser interpretados à luz das
questões de investigação da Secção~\ref{sec:intro:questions} e comparados
com os trabalhos apresentados no Capítulo~\ref{chap:review}. A discussão
deve explicar porque é que os resultados são os que são, e não apenas
repeti-los.

\end{otherlanguage}
